\chapter{Importy, \texttt{@include} i skladanie projektu z wielu plikow}

\section{Dlaczego w ogole dzielic projekt na pliki}
Na poczatku jeden plik wydaje sie wygodny. Jednak wraz z rozwojem projektu zaczyna byc
za duzy: miesza metadata, logike, strony, dane pomocnicze i drobne narzedzia. Wtedy
naturalnym krokiem jest podzial na mniejsze elementy.

MoonChunk obsluguje taki podzial przez:

\begin{itemize}
  \item \mc{import},
  \item \mc{export chunk},
  \item \mc{@include},
  \item jawne wywolania \mc{moon(...)}.
\end{itemize}

\section{Importy musza stac na samej gorze pliku}
To jedna z zasad organizacyjnych MoonChunk. Importy powinny pojawic sie przed deklaracjami
chunkow. Dzięki temu od razu widac zaleznosci danego pliku.

\begin{lstlisting}[style=moonchunk,caption={Importy na samej gorze pliku}]
import { Shared } from "./shared.mncnk";
import * as Pages from "./pages.mncnk";

chunk "Main" {
  output: "./dist";
};
\end{lstlisting}

To nie jest tylko kwestia estetyki. Taki uklad wspiera przewidywalne parsowanie i czytelna
organizacje kodu.

\section{Eksportowany chunk}
Aby chunk mogl byc importowany z innego pliku, powinien byc oznaczony slowem \mc{export}.

\begin{lstlisting}[style=moonchunk,caption={Chunk eksportowany do innych plikow}]
export chunk "Metadata" {
  title: "MoonChunk";
};
\end{lstlisting}

Taki zapis komunikuje jasno: ten chunk jest elementem publicznym, z ktorego inne pliki
moga korzystac.

\section{Import nazwany}
Import nazwany pozwala pobrac konkretne chunki po nazwie.

\begin{lstlisting}[style=moonchunk,caption={Import nazwany}]
import { Metadata } from "./metadata.mncnk";
\end{lstlisting}

To dobry wybor wtedy, gdy potrzebujesz jednego lub kilku konkretnych elementow.

\section{Import przestrzeni nazw}
Druga forma to import calej przestrzeni nazw pod aliasem.

\begin{lstlisting}[style=moonchunk,caption={Import przestrzeni nazw}]
import * as Fragments from "./fragments.mncnk";
\end{lstlisting}

Potem mozesz odwolywac sie do elementow przez kropke, np. \mc{Fragments.PageHeader}.
To szczegolnie wygodne, gdy z jednego pliku chcesz udostepnic wieksza kolekcje chunkow.

\section{Samo \mc{import} niczego jeszcze nie uruchamia}
To bardzo wazna zasada: import jedynie sprawia, ze chunki staja sie dostepne. Nie oznacza
automatycznego wykonania.

Jezeli zaimportujesz plik, ale nie uzyjesz zadnego z jego chunkow przez \mc{@include}
albo \mc{moon(...)}, to ten kod nie zostanie uruchomiony tylko dlatego, ze istnieje.
To zachowanie celowo chroni projekt przed ukrytymi efektami ubocznymi.

\section{Do czego sluzy \mc{@include}}
\mc{@include} wlacza zawartosc jednego chunku do drugiego. To bardzo wygodne przy
skladaniu wspolnych metadata, wspolnej logiki lub fragmentow konfiguracyjnych.

\begin{lstlisting}[style=moonchunk,caption={Dolaczenie wspolnego chunku}]
import { Metadata } from "./metadata.mncnk";

chunk "Main" {
  @include Metadata;
  output: "./dist";
};
\end{lstlisting}

Mozesz o tym myslec jak o swiadomym powiedzeniu: \emph{do tego chunku dolacz tez ten
inny, wspolny zestaw instrukcji}.

\section{\mc{@include} powinno pojawiac sie na poczatku chunku}
MoonChunk narzuca uporzadkowany styl: instrukcje \mc{@include} powinny stac na samym
poczatku ciala chunku. To sprawia, ze od razu widac, od czego dany chunk zalezy.

Taki uklad jest czytelny i wart zachowania nawet wtedy, gdy projekt rośnie.

\section{Import nazwany a import przez namespace}
Roznica praktyczna jest prosta.

\subsection*{Import nazwany}
\begin{lstlisting}[style=moonchunk,caption={Uzycie importu nazwanego}]
import { Metadata } from "./metadata.mncnk";

chunk "Main" {
  @include Metadata;
};
\end{lstlisting}

\subsection*{Import przestrzeni nazw}
\begin{lstlisting}[style=moonchunk,caption={Uzycie importu namespace}]
import * as Fragments from "./fragments.mncnk";

chunk "Main" {
  @include Fragments.Metadata;
};
\end{lstlisting}

Pierwsza forma jest krotsza. Druga jest bardziej jawna i wygodna, gdy importujesz
wiecej rzeczy z jednego miejsca.

\section{Lancuch importow}
MoonChunk obsluguje projekty, w ktorych jeden plik importuje drugi, a ten kolejny.
Warto jednak zachowac porzadek i unikac zbednej glebokosci. Im prostszy graf zaleznosci,
tym latwiej zrozumiec projekt.

\section{Cykliczne zaleznosci}
Jesli dwa pliki zaczelyby importowac sie wzajemnie w niebezpieczny sposob, MoonChunk
powinien zglosic blad zwiazany z petla importow. To dobra ochrona przed konfiguracja,
ktora bylaby bardzo trudna do przewidzenia dla uzytkownika.

\section{Punkt wejscia nadal musi byc jawny}
Nawet jesli projekt sklada sie z wielu plikow i wielu chunkow, nadal ostatecznie ktos
musi wybrac, co ma zostac wykonane. Robi to \mc{moon(...)}.

\begin{lstlisting}[style=moonchunk,caption={Jawne wywolania punktow wejscia}]
moon(Main);
moon(Secondary);
\end{lstlisting}

Dzieki temu importy i \mc{@include} odpowiadaja za \emph{dostepnosc i kompozycje},
a \mc{moon(...)} odpowiada za \emph{faktyczne uruchomienie}.

\section{Typowy wzorzec modularnego projektu}
Bardzo praktyczny uklad wyglada tak:

\begin{lstlisting}[style=basecode,caption={Przykladowy projekt wieloplikowy}]
content/
  site.mncnk
  metadata.mncnk
  sections.mncnk
  blog.mncnk
  data/
    posts.json
\end{lstlisting}

Możliwe role:

\begin{itemize}
  \item \mc{site.mncnk} -- glowny plik z punktami \mc{moon(...)},
  \item \mc{metadata.mncnk} -- pola wspolne dla calej witryny,
  \item \mc{sections.mncnk} -- dodatkowe chunki do dolaczania,
  \item \mc{blog.mncnk} -- logika stron blogowych,
  \item \mc{posts.json} -- dane wczytywane przez \mc{data(...)}.
\end{itemize}

\section{Najczestsze bledy przy importach}
\begin{itemize}
  \item import w srodku pliku zamiast na gorze,
  \item probowanie importu chunku bez \mc{export},
  \item literowka w nazwie importowanego chunku,
  \item probowanie uzycia aliasu namespace, ktory nie istnieje,
  \item oczekiwanie, ze sam import uruchomi generowanie strony.
\end{itemize}

\section{Jak myslec o \mc{@include}}
Najprostsza intuicja jest taka: import otwiera dostep do chunku, a \mc{@include}
mowi, zeby ten chunk realnie dolaczyc do aktualnego miejsca. To rozdzielenie jest
bardzo zdrowe architektonicznie, bo oddziela \emph{zaleznosc} od \emph{wykonania}.
